\setcounter{chapter}{-1}
\chapter{预备知识}

为便于使用，本章汇集了贯穿全书将要用到的一些结果和记号。学生起初可以快速浏览本章，随后在正文中遇到相应概念时，再更仔细地阅读各节。

\section{基础}

集合论的基础知识：集合、$\cap$、$\cup$、$\in$ 等，应为读者所熟悉。对于给定集合 $A$ 的子集，我们将使用记号
\[
  B=\{a\in A\mid \cdots\text{（关于 }a\text{ 的条件）}\cdots\}.
\]
集合 $A$ 的\emph{阶}或\emph{基数}记为 $|A|$。若 $A$ 是有限集，则 $A$ 的阶就是 $A$ 中元素的个数。

重要的是要理解如何检验某个特定的 $x\in A$ 是否属于 $A$ 的一个子集 $B$（参见练习 1--4）。两个集合 $A$ 与 $B$ 的\emph{笛卡儿积}是由来自 $A$ 和 $B$ 的有序对组成的集合
\[
  A\times B=\{(a,b)\mid a\in A,\ b\in B\}.
\]

对于一些常见的数集，我们将使用如下记号：
\begin{enumerate}[label=(\arabic*)]
  \item $\ZZ=\{0,\pm1,\pm2,\pm3,\ldots\}$ 表示整数（字母 $Z$ 来自德语中表示“数”的词 ``Zahlen''）。
  \item $\QQ=\{a/b\mid a,b\in\ZZ,\ b\neq0\}$ 表示有理数。
  \item $\RR=\{\text{所有十进制展开 }\pm d_1d_2\cdots d_n.a_1a_2a_3\cdots\}$ 表示实数。
  \item $\CC=\{a+bi\mid a,b\in\RR,\ i^2=-1\}$ 表示复数。
  \item $\ZZ^+$、$\QQ^+$ 和 $\RR^+$ 分别表示 $\ZZ$、$\QQ$ 和 $\RR$ 中的正的（非零）元素。
\end{enumerate}

我们将使用记号 $f:A\to B$ 或 $A\xrightarrow{f}B$ 表示从 $A$ 到 $B$ 的函数 $f$，$f$ 在 $a$ 处的值记为 $f(a)$（即所有函数都写在自变量的左侧）。我们交替使用“函数”和“映射”这两个词。集合 $A$ 称为 $f$ 的\emph{定义域}，$B$ 称为 $f$ 的\emph{陪域}。记号 $f:a\mapsto b$，或者在 $f$ 已经明确时写作 $a\mapsto b$，表示 $f(a)=b$，也就是说，函数是在元素上被指定的。

如果函数 $f$ 不是在元素上给定的，那么通常必须检验 $f$ 是\emph{良定}的，即它是无歧义地确定的。例如，若集合 $A$ 是两个子集 $A_1$ 与 $A_2$ 的并，则可以尝试指定一个从 $A$ 到集合 $\{0,1\}$ 的函数：规定 $A_1$ 中的所有元素都映到 $0$，$A_2$ 中的所有元素都映到 $1$。除非 $A_1$ 与 $A_2$ 有公共元素，否则这就无歧义地定义了 $f$（若它们有公共元素，则不清楚这些元素应映到 $0$ 还是 $1$）。因此，检验这个 $f$ 是否良定，就等同于检验 $A_1$ 与 $A_2$ 没有交集。

集合
\[
  f(A)=\{b\in B\mid b=f(a),\ \text{对某个 }a\in A\}
\]
是 $B$ 的一个子集，称为 $f$ 的\emph{值域}或\emph{像}（也称为 $A$ 在 $f$ 下的像）。对 $B$ 的每个子集 $C$，集合
\[
  f^{-1}(C)=\{a\in A\mid f(a)\in C\}
\]
由 $A$ 中在 $f$ 下映入 $C$ 的元素组成，称为 $C$ 在 $f$ 下的\emph{原像}或\emph{逆像}。对每个 $b\in B$，$\{b\}$ 在 $f$ 下的原像称为 $f$ 在 $b$ 上方的\emph{纤维}。注意，一般而言 $f^{-1}$ 并不是函数，并且 $f$ 的纤维通常含有许多元素，因为可能有许多 $A$ 中的元素映到同一个元素 $b$。

若 $f:A\to B$ 且 $g:B\to C$，则复合映射 $g\circ f:A\to C$ 定义为
\[
  (g\circ f)(a)=g(f(a)).
\]

设 $f:A\to B$。
\begin{enumerate}[label=(\arabic*)]
  \item 若只要 $a_1\neq a_2$ 就有 $f(a_1)\neq f(a_2)$，则称 $f$ 是\emph{单射}的，或称 $f$ 是一个\emph{单射}。
  \item 若对所有 $b\in B$，存在某个 $a\in A$ 使得 $f(a)=b$，即 $f$ 的像是整个 $B$，则称 $f$ 是\emph{满射}的，或称 $f$ 是一个\emph{满射}。注意，因为函数总是映到它的值域上（按定义），所以必须指定陪域 $B$，满射性的问题才有意义。
  \item 若 $f$ 既是单射又是满射，则称 $f$ 是\emph{双射}的，或称 $f$ 是一个\emph{双射}。若存在这样一个从 $A$ 到 $B$ 的双射 $f$，则称 $A$ 与 $B$ 是\emph{双射对应}的。
  \item 若存在函数 $g:B\to A$，使得 $g\circ f:A\to A$ 是 $A$ 上的恒等映射，即对所有 $a\in A$ 都有 $(g\circ f)(a)=a$，则称 $f$ 有一个\emph{左逆}。
  \item 若存在函数 $h:B\to A$，使得 $f\circ h:B\to B$ 是 $B$ 上的恒等映射，则称 $f$ 有一个\emph{右逆}。
\end{enumerate}

\medskip
\noindent\textbf{命题 1.} 设 $f:A\to B$。
\begin{enumerate}[label=(\arabic*)]
  \item 映射 $f$ 是单射，当且仅当 $f$ 有左逆。
  \item 映射 $f$ 是满射，当且仅当 $f$ 有右逆。
  \item 映射 $f$ 是双射，当且仅当存在 $g:B\to A$，使得 $f\circ g$ 是 $B$ 上的恒等映射，并且 $g\circ f$ 是 $A$ 上的恒等映射。
  \item 若 $A$ 与 $B$ 是元素个数相同的有限集（即 $|A|=|B|$），则 $f:A\to B$ 是双射，当且仅当 $f$ 是单射，当且仅当 $f$ 是满射。
\end{enumerate}
\noindent\textit{证明：} 练习。

在上面命题第 (3) 部分的情形中，映射 $g$ 必然唯一；我们称 $g$ 为 $f$ 的\emph{双侧逆}（或简称为\emph{逆}）。

集合 $A$ 的一个\emph{置换}就是从 $A$ 到自身的一个双射。

若 $A\subseteq B$ 且 $f:B\to C$，则 $f$ 到 $A$ 的限制记为 $f|_A$。当所考虑的定义域已明确时，我们偶尔仍把 $f|_A$ 简写为 $f$，尽管形式上它们是不同的函数（因为它们的定义域不同）。

若 $A\subseteq B$，$g:A\to C$，并且存在函数 $f:B\to C$ 使得 $f|_A=g$，则称 $f$ 是 $g$ 到 $B$ 的一个\emph{延拓}（这样的映射 $f$ 不一定存在，也不一定唯一）。

设 $A$ 是非空集合。
\begin{enumerate}[label=(\arabic*)]
  \item 集合 $A$ 上的一个\emph{二元关系}是 $A\times A$ 的一个子集 $R$，若 $(a,b)\in R$，则写作 $a\sim b$。
  \item 称 $A$ 上的关系 $\sim$ 为：
  \begin{enumerate}[label=(\alph*)]
    \item \emph{自反的}，若对所有 $a\in A$，有 $a\sim a$；
    \item \emph{对称的}，若对所有 $a,b\in A$，$a\sim b$ 蕴含 $b\sim a$；
    \item \emph{传递的}，若对所有 $a,b,c\in A$，$a\sim b$ 且 $b\sim c$ 蕴含 $a\sim c$。
  \end{enumerate}
  若一个关系是自反的、对称的并且传递的，则称它是一个\emph{等价关系}。
  \item 若 $\sim$ 在 $A$ 上定义了一个等价关系，则 $a\in A$ 的\emph{等价类}定义为 $\{x\in A\mid x\sim a\}$。$a$ 的等价类中的元素称为与 $a$ 等价。若 $C$ 是一个等价类，则 $C$ 中的任一元素都称为类 $C$ 的一个\emph{代表元}。
  \item $A$ 的一个\emph{划分}是由 $A$ 的非空子集组成的任意集合 $\{A_i\mid i\in I\}$（其中 $I$ 是某个指标集），满足
  \begin{enumerate}[label=(\alph*)]
    \item $A=\bigcup_{i\in I}A_i$，并且
    \item 对所有满足 $i\neq j$ 的 $i,j\in I$，有 $A_i\cap A_j=\varnothing$，
  \end{enumerate}
  也就是说，$A$ 是该划分中各集合的不交并。
\end{enumerate}

$A$ 上的等价关系与 $A$ 的划分是同一回事：

\medskip
\noindent\textbf{命题 2.} 设 $A$ 是非空集合。
\begin{enumerate}[label=(\arabic*)]
  \item 若 $\sim$ 在 $A$ 上定义了一个等价关系，则 $\sim$ 的等价类全体构成 $A$ 的一个划分。
  \item 若 $\{A_i\mid i\in I\}$ 是 $A$ 的一个划分，则 $A$ 上存在一个等价关系，其等价类恰好是这些集合 $A_i$，$i\in I$。
\end{enumerate}
\noindent\textit{证明：} 略。

最后，我们假定读者熟悉归纳法证明。

\subsection*{练习}

在练习 1 到 4 中，令 $\mathcal{A}$ 为实数域上所有 $2\times2$ 矩阵组成的集合。回忆矩阵乘法定义为
\[
  \begin{pmatrix}a&b\\ c&d\end{pmatrix}
  \begin{pmatrix}p&q\\ r&s\end{pmatrix}
  =
  \begin{pmatrix}ap+br&aq+bs\\ cp+dr&cq+ds\end{pmatrix}.
\]
令
\[
  M=\begin{pmatrix}1&1\\0&1\end{pmatrix}
\]
并令
\[
  \mathcal{B}=\{X\in\mathcal{A}\mid MX=XM\}.
\]
\begin{enumerate}[label=\arabic*.]
  \item 判断下列 $\mathcal{A}$ 中的元素哪些属于 $\mathcal{B}$：
  \[
  \begin{pmatrix}1&1\\0&1\end{pmatrix},\quad
  \begin{pmatrix}1&1\\1&1\end{pmatrix},\quad
  \begin{pmatrix}0&0\\0&0\end{pmatrix},\quad
  \begin{pmatrix}1&1\\1&0\end{pmatrix},\quad
  \begin{pmatrix}1&0\\0&1\end{pmatrix},\quad
  \begin{pmatrix}0&1\\1&0\end{pmatrix}.
  \]
  \item 证明：若 $P,Q\in\mathcal{B}$，则 $P+Q\in\mathcal{B}$（其中 $+$ 表示通常的矩阵加法）。
  \item 证明：若 $P,Q\in\mathcal{B}$，则 $P\cdot Q\in\mathcal{B}$（其中 $\cdot$ 表示通常的矩阵乘法）。
  \item 求出关于 $p,q,r,s$ 的条件，它们精确地刻画何时 $\begin{pmatrix}p&q\\ r&s\end{pmatrix}\in\mathcal{B}$。
  \item 判断下列函数 $f$ 是否良定：
  \begin{enumerate}[label=(\alph*)]
    \item $f:\QQ\to\ZZ$，定义为 $f(a/b)=a$。
    \item $f:\QQ\to\QQ$，定义为 $f(a/b)=a^2/b^2$。
  \end{enumerate}
  \item 判断函数 $f:\RR^+\to\ZZ$ 是否良定，其中 $f$ 把实数 $r$ 映到 $r$ 的某个十进制展开中小数点右边的第一位数字。
  \item 设 $f:A\to B$ 是集合之间的满射。证明关系
  \[
    a\sim b\quad\text{当且仅当}\quad f(a)=f(b)
  \]
  是一个等价关系，其等价类就是 $f$ 的纤维。
\end{enumerate}

\section{整数的性质}

整数 $\ZZ$ 的以下性质（其中许多是初等算术中熟悉的性质）将在第 8 章的环论中置于更一般的背景下证明；不过在第一部分中有必要使用它们（当然，环论中对这些性质的证明并不依赖群论）。
\begin{enumerate}[label=(\arabic*)]
  \item （$\ZZ$ 的良序性）若 $A$ 是 $\ZZ^+$ 的任一非空子集，则存在某个元素 $m\in A$，使得对所有 $a\in A$ 都有 $m\leq a$（$m$ 称为 $A$ 的一个\emph{极小元}）。
  \item 若 $a,b\in\ZZ$ 且 $a\neq0$，如果存在元素 $c\in\ZZ$ 使得 $b=ac$，则称 $a$\emph{整除} $b$。此时记为 $a\mid b$；若 $a$ 不整除 $b$，则记为 $a\nmid b$。
  \item 若 $a,b\in\ZZ-\{0\}$，则存在唯一的正整数 $d$，称为 $a$ 与 $b$ 的\emph{最大公因子}（或 g.c.d.），满足：
  \begin{enumerate}[label=(\alph*)]
    \item $d\mid a$ 且 $d\mid b$（因此 $d$ 是 $a$ 与 $b$ 的一个公因子），并且
    \item 若 $e\mid a$ 且 $e\mid b$，则 $e\mid d$（因此 $d$ 是这样的最大因子）。
  \end{enumerate}
  $a$ 与 $b$ 的最大公因子记为 $(a,b)$。若 $(a,b)=1$，则称 $a$ 与 $b$\emph{互素}。
  \item 若 $a,b\in\ZZ-\{0\}$，则存在唯一的正整数 $l$，称为 $a$ 与 $b$ 的\emph{最小公倍数}（或 l.c.m.），满足：
  \begin{enumerate}[label=(\alph*)]
    \item $a\mid l$ 且 $b\mid l$（因此 $l$ 是 $a$ 与 $b$ 的一个公倍数），并且
    \item 若 $a\mid m$ 且 $b\mid m$，则 $l\mid m$（因此 $l$ 是这样的最小倍数）。
  \end{enumerate}
  两个整数 $a$ 与 $b$ 的最大公因子 $d$ 与最小公倍数 $l$ 之间的联系由 $dl=ab$ 给出。
  \item \emph{带余除法}：若 $a,b\in\ZZ-\{0\}$，则存在唯一的 $q,r\in\ZZ$，使得
  \[
    a=qb+r\quad\text{且}\quad 0\leq r<|b|,
  \]
  其中 $q$ 是商，$r$ 是余数。这就是初等算术中熟悉的通常“长除法”。
  \item \emph{欧几里得算法}是一个重要过程，它通过反复使用带余除法来求出两个整数 $a$ 与 $b$ 的最大公因子：若 $a,b\in\ZZ-\{0\}$，则可得到一列商与余数
  \begin{alignat*}{2}
    a     &= q_0b+r_0,       &\qquad &(0)\\
    b     &= q_1r_0+r_1,     &&(1)\\
    r_0   &= q_2r_1+r_2,     &&(2)\\
    r_1   &= q_3r_2+r_3,     &&(3)\\
          &\ \vdots          &&\\
    r_{n-2}&=q_nr_{n-1}+r_n, &&(n)\\
    r_{n-1}&=q_{n+1}r_n,     &&(n+1)
  \end{alignat*}
  其中 $r_n$ 是最后一个非零余数。这样的 $r_n$ 必然存在，因为若余数非零，则 $|b|>|r_0|>|r_1|>\cdots>|r_n|$ 是严格正整数的递减序列，而这样的序列不可能无限继续。于是 $r_n$ 就是 $a$ 与 $b$ 的最大公因子 $(a,b)$。
\end{enumerate}

\paragraph{例}
设 $a=57970$，$b=10353$。应用欧几里得算法得到：
\begin{align*}
  57970&=(5)10353+6205,\\
  10353&=(1)6205+4148,\\
  6205&=(1)4148+2057,\\
  4148&=(2)2057+34,\\
  2057&=(60)34+17,\\
  34&=(2)17,
\end{align*}
这说明 $(57970,10353)=17$。

\begin{enumerate}[label=(\arabic*),resume]
  \item 欧几里得算法的一个后果将被我们经常使用：若 $a,b\in\ZZ-\{0\}$，则存在 $x,y\in\ZZ$，使得
  \[
    (a,b)=ax+by,
  \]
  也就是说，$a$ 与 $b$ 的最大公因子是 $a$ 与 $b$ 的一个 $\ZZ$-线性组合。这可由在欧几里得算法中递归地把元素 $r_n$ 写成前面余数的组合而得出（具体地说，利用上面的方程 $(n)$ 解出 $r_n=r_{n-2}-q_nr_{n-1}$，从而用余数 $r_{n-1}$ 与 $r_{n-2}$ 表示 $r_n$；再利用方程 $(n-1)$ 用 $r_{n-2}$ 与 $r_{n-3}$ 表示 $r_n$，依此类推，最终用 $a$ 与 $b$ 表示 $r_n$）。
\end{enumerate}

\paragraph{例}
设 $a=57970$，$b=10353$，它们的最大公因子上面已经算出为 $17$。对这两个整数应用欧几里得算法时，从第五个方程（倒数第二个方程）解出它们的最大公因子：
\[
  17=2057-(60)34.
\]
第四个方程又给出 $34=4148-(2)2057$，所以把前一个余数 $34$ 的这个表达式代入，得到
\[
  17=2057-(60)[4148-(2)2057],
\]
即
\[
  17=(121)2057-(60)4148.
\]
由第三个方程解出 $2057$ 并代入，得
\[
  17=(121)[6205-(1)4148]-(60)4148=(121)6205-(181)4148.
\]
再用第二个方程解出 $4148$，然后用第一个方程解出 $6205$，最终得到
\[
  17=(302)57970-(1691)10353,
\]
这一点很容易直接检验。因此，在这个例子中，对于最大公因子的方程 $ax+by=(a,b)$，有解 $x=302$ 与 $y=-1691$。注意，单靠猜测找到这个关系的可能性相当小。

上面 (7) 中的整数 $x$ 与 $y$ 并不唯一。例如，在 $a=57970$、$b=10353$ 的例子中，我们求得一个解为 $x=302$、$y=-1691$；同时也比较容易检验 $x=-307$、$y=1719$ 也满足 $57970x+10353y=17$。$x$ 与 $y$ 的一般解是已知的（参见下面的练习以及第 8 章）。

\begin{enumerate}[label=(\arabic*),resume]
  \item $\ZZ^+$ 的一个元素 $p$ 称为\emph{素数}，若 $p>1$ 且 $p$ 的正因子只有 $1$ 和 $p$（起初，“素数”一词只指正整数）。大于 $1$ 但不是素数的整数称为\emph{合数}。例如，$2,3,5,7,11,13,17,19,\ldots$ 是素数，而 $4,6,8,9,10,12,14,15,16,18,\ldots$ 是合数。

  素数的一个重要性质（事实上可以用来定义素数，参见练习 3）如下：若 $p$ 是素数，并且对某些 $a,b\in\ZZ$ 有 $p\mid ab$，则 $p\mid a$ 或 $p\mid b$。

  \item \emph{算术基本定理}说：若 $n\in\ZZ$ 且 $n>1$，则 $n$ 可以唯一地分解为素数的乘积；也就是说，存在互不相同的素数 $p_1,p_2,\ldots,p_s$ 和正整数 $\alpha_1,\alpha_2,\ldots,\alpha_s$，使得
  \[
    n=p_1^{\alpha_1}p_2^{\alpha_2}\cdots p_s^{\alpha_s}.
  \]
  这个分解的唯一性是指：若 $q_1,q_2,\ldots,q_t$ 是任意互不相同的素数，$\beta_1,\beta_2,\ldots,\beta_t$ 是正整数，并且
  \[
    n=q_1^{\beta_1}q_2^{\beta_2}\cdots q_t^{\beta_t},
  \]
  则 $s=t$，并且如果把两组素数都按递增顺序排列，则 $q_i=p_i$ 且 $\alpha_i=\beta_i$，$1\leq i\leq s$。例如
  \[
    n=1852423848=2^3 3^2 11^2 19^3 31,
  \]
  而这个素数乘积分解是唯一的。

  假设正整数 $a$ 与 $b$ 表示为素数幂的乘积：
  \[
    a=p_1^{\alpha_1}p_2^{\alpha_2}\cdots p_s^{\alpha_s},
    \qquad
    b=p_1^{\beta_1}p_2^{\beta_2}\cdots p_s^{\beta_s},
  \]
  其中 $p_1,p_2,\ldots,p_s$ 互不相同，并且指数都 $\geq0$（这里允许指数为 $0$，以便两个乘积取在同一组素数上；若某个素数实际上不是因子，则其指数为 $0$）。于是 $a$ 与 $b$ 的最大公因子为
  \[
    (a,b)=p_1^{\min(\alpha_1,\beta_1)}p_2^{\min(\alpha_2,\beta_2)}\cdots p_s^{\min(\alpha_s,\beta_s)}
  \]
  （最小公倍数则改为取 $\alpha_i$ 与 $\beta_i$ 的最大值，而不是最小值）。
\end{enumerate}

\paragraph{例}
在上面的例子中，$a=57970$ 与 $b=10353$ 可分解为
\[
  a=2\cdot5\cdot11\cdot17\cdot31,
  \qquad
  b=3\cdot7\cdot17\cdot29,
\]
由此可以立即得出它们的最大公因子为 $17$。不过要注意，对大整数而言，确定其素因数分解是极其困难的（事实上，目前使用的若干常见编码正是基于这种困难），因此这通常不是求最大公因子的有效方法。欧几里得算法无需 $a$ 与 $b$ 的素因数分解，就能相当快速地求出最大公因子。

\begin{enumerate}[label=(\arabic*),resume]
  \item \emph{欧拉 $\varphi$ 函数}定义如下：对 $n\in\ZZ^+$，令 $\varphi(n)$ 为满足 $a\leq n$ 且 $a$ 与 $n$ 互素，即 $(a,n)=1$ 的正整数 $a$ 的个数。例如，$\varphi(12)=4$，因为 $1,5,7,11$ 是不超过 $12$ 且与 $12$ 没有公共因子的唯一正整数。同样，$\varphi(1)=1$，$\varphi(2)=1$，$\varphi(3)=2$，$\varphi(4)=2$，$\varphi(5)=4$，$\varphi(6)=2$，等等。对素数 $p$，有 $\varphi(p)=p-1$；更一般地，对所有 $a\geq1$，有公式
  \[
    \varphi(p^a)=p^a-p^{a-1}=p^{a-1}(p-1).
  \]
  函数 $\varphi$ 在如下意义下是\emph{乘性的}：
  \[
    \varphi(ab)=\varphi(a)\varphi(b)\qquad\text{若 }(a,b)=1
  \]
  （注意这里 $a$ 与 $b$ 互素是重要的）。结合上面的公式，可得到 $\varphi$ 的一般取值公式：若 $n=p_1^{\alpha_1}p_2^{\alpha_2}\cdots p_s^{\alpha_s}$，则
  \begin{align*}
    \varphi(n)
      &=\varphi(p_1^{\alpha_1})\varphi(p_2^{\alpha_2})\cdots\varphi(p_s^{\alpha_s})\\
      &=p_1^{\alpha_1-1}(p_1-1)p_2^{\alpha_2-1}(p_2-1)\cdots p_s^{\alpha_s-1}(p_s-1).
  \end{align*}
  例如，$\varphi(12)=\varphi(2^2)\varphi(3)=2^1(2-1)3^0(3-1)=4$。读者应注意，在全书中我们会用字母 $\varphi$ 表示许多不同的函数，因此当我们要用这个字母表示欧拉函数时，会明确说明。
\end{enumerate}

\subsection*{练习}

\begin{enumerate}[label=\arabic*.]
  \item 对下列每一对整数 $a$ 与 $b$，求出它们的最大公因子、最小公倍数，并把它们的最大公因子写成 $ax+by$ 的形式，其中 $x,y$ 是整数。
  \begin{enumerate}[label=(\alph*)]
    \item $a=20,\ b=13$。
    \item $a=69,\ b=372$。
    \item $a=792,\ b=275$。
    \item $a=11391,\ b=5673$。
    \item $a=1761,\ b=1567$。
    \item $a=507885,\ b=60808$。
  \end{enumerate}
  \item 证明：若整数 $k$ 整除整数 $a$ 与 $b$，则对任意一对整数 $s$ 与 $t$，$k$ 整除 $as+bt$。
  \item 证明：若 $n$ 是合数，则存在整数 $a$ 与 $b$，使得 $n$ 整除 $ab$，但 $n$ 不整除 $a$ 也不整除 $b$。
  \item 设 $a,b,N$ 是固定整数，其中 $a$ 与 $b$ 非零；令 $d=(a,b)$ 为 $a$ 与 $b$ 的最大公因子。假设 $x_0$ 与 $y_0$ 是 $ax+by=N$ 的一组特解（即 $ax_0+by_0=N$）。证明：对任意整数 $t$，整数
  \[
    x=x_0+\frac{b}{d}t,\qquad y=y_0-\frac{a}{d}t
  \]
  也是 $ax+by=N$ 的解（这事实上给出了一般解）。
  \item 求出每个整数 $n\leq30$ 的 $\varphi(n)$ 值，其中 $\varphi$ 表示欧拉 $\varphi$ 函数。
  \item 用归纳法证明 $\ZZ$ 的良序性，并证明极小元唯一。
  \item 若 $p$ 是素数，证明不存在非零整数 $a$ 与 $b$ 使得 $a^2=pb^2$（即 $\sqrt{p}$ 不是有理数）。
  \item 设 $p$ 是素数，$n\in\ZZ^+$。求出整除
  \[
    n!=n(n-1)(n-2)\cdots 2\cdot1
  \]
  的最大 $p$ 的幂的公式（它涉及取整函数）。
  \item 编写一个计算机程序，求两个整数 $a$ 与 $b$ 的最大公因子 $(a,b)$，并把 $(a,b)$ 表示为 $ax+by$，其中 $x,y$ 是整数。
  \item 证明：对任意给定的正整数 $N$，只有有限多个整数 $n$ 满足 $\varphi(n)=N$，其中 $\varphi$ 表示欧拉 $\varphi$ 函数。特别地，由此推出当 $n$ 趋于无穷时，$\varphi(n)$ 也趋于无穷。
  \item 证明：若 $d$ 整除 $n$，则 $\varphi(d)$ 整除 $\varphi(n)$，其中 $\varphi$ 表示欧拉 $\varphi$ 函数。
\end{enumerate}

\section{$\ZZ/n\ZZ$：模 $n$ 的整数}

令 $n$ 是一个固定的正整数。在 $\ZZ$ 上定义关系：
\[
  a\sim b\quad\text{当且仅当}\quad n\mid(b-a).
\]
显然，对任意整数 $a$，有 $a\sim a$；并且对任意整数 $a,b$，$a\sim b$ 蕴含 $b\sim a$，所以这个关系显然是自反且对称的。若 $a\sim b$ 且 $b\sim c$，则 $n$ 整除 $a-b$，并且 $n$ 整除 $b-c$，于是 $n$ 也整除这两个整数之和，即 $n$ 整除
\[
  (a-b)+(b-c)=a-c,
\]
所以 $a\sim c$，从而这个关系是传递的。因此这是一个等价关系。若 $a\sim b$，则写作 $a\equiv b\pmod n$（读作：$a$ 模 $n$ 同余于 $b$）。对任意 $k\in\ZZ$，我们把 $a$ 的等价类记为 $\overline{a}$；这称为 $a$ 模 $n$ 的\emph{同余类}或\emph{剩余类}，它由与 $a$ 相差 $n$ 的整数倍的所有整数组成，即
\[
  \overline{a}=\{a+kn\mid k\in\ZZ\}
  =\{a,\ a\pm n,\ a\pm2n,\ a\pm3n,\ldots\}.
\]

模 $n$ 恰有 $n$ 个不同的等价类，即
\[
  \overline{0},\overline{1},\overline{2},\ldots,\overline{n-1},
\]
它们由除以 $n$ 后可能出现的余数决定，并且这些剩余类划分整数集 $\ZZ$。在这个等价关系下的等价类集合记为 $\ZZ/n\ZZ$，称为\emph{模 $n$ 的整数}（或\emph{模 $n$ 整数}）。当我们讨论商群和商环时，这个记号的动机会变得更加清楚。注意，对于不同的 $n$，等价关系和等价类是不同的；因此在使用横线记号以前，我们总是会先固定 $n$。求某个整数 $a$ 模 $n$ 的等价类这一过程常称为把 $a$\emph{化简模 $n$}。这个术语也常指求出与 $a$ 模 $n$ 同余的最小非负整数（$a$ 模 $n$ 的\emph{最小剩余}）。

我们可以为 $\ZZ/n\ZZ$ 的元素定义加法和乘法，从而定义\emph{模算术}如下：对 $\overline{a},\overline{b}\in\ZZ/n\ZZ$，定义它们的和与积为
\[
  \overline{a}+\overline{b}=\overline{a+b},
  \qquad
  \overline{a}\cdot\overline{b}=\overline{ab}.
\]
这意味着如下做法：给定 $\ZZ/n\ZZ$ 中的任意两个元素 $\overline{a}$ 与 $\overline{b}$，要计算它们的和（分别地，积），取类 $\overline{a}$ 中任意一个代表整数 $a$，以及类 $\overline{b}$ 中任意一个代表整数 $b$，在 $\ZZ$ 中照常把整数 $a$ 与 $b$ 相加（分别地，相乘），然后取包含结果的等价类。下面的定理 3 断言这是良定的，即它不依赖于为 $\ZZ/n\ZZ$ 中的元素 $\overline{a}$ 与 $\overline{b}$ 所选择的代表元。

\paragraph{例}
设 $n=12$，考虑 $\ZZ/12\ZZ$，它由十二个剩余类
\[
  \overline{0},\overline{1},\overline{2},\ldots,\overline{11}
\]
组成，这些类由整数除以 $12$ 后的十二种可能余数决定。例如，剩余类 $\overline{5}$ 中的元素，就是除以 $12$ 时余数为 $5$ 的那些整数（即模 $12$ 同余于 $5$ 的整数）。任意一个模 $12$ 同余于 $5$ 的整数（如 $5,17,29,\ldots$ 或 $-7,-19,\ldots$）都可作为剩余类 $\overline{5}$ 的一个代表元。注意，$\ZZ/12\ZZ$ 由上述十二个元素组成（而 $\ZZ/12\ZZ$ 的这些元素中，每一个本身又包含无限多个通常意义下的整数）。

现在设 $\overline{a}=\overline{5}$ 且 $\overline{b}=\overline{8}$。$\overline{a}$ 最显然的代表元是整数 $5$，同样 $\overline{b}$ 最显然的代表元是 $8$。使用这些代表元，得到
\[
  \overline{5}+\overline{8}=\overline{13}=\overline{1},
\]
因为 $13$ 与 $1$ 属于模 $n=12$ 的同一个类。若改用代表元 $17$ 代表 $\overline{a}$（注意 $5$ 与 $17$ 的确属于模 $12$ 的同一个剩余类），并用代表元 $-28$ 代表 $\overline{b}$，则得到
\[
  \overline{5}+\overline{8}=\overline{17-28}=\overline{-11}=\overline{1},
\]
如前所述，结果不依赖于所选择的代表元。这两个类的积为
\[
  \overline{a}\cdot\overline{b}=\overline{5}\cdot\overline{8}=\overline{40}=\overline{4},
\]
同样与所选择的代表元无关。

\medskip
\noindent\textbf{定理 3.} 上面定义的 $\ZZ/n\ZZ$ 上的加法与乘法运算都是良定的，也就是说，它们不依赖于所涉及各类代表元的选择。更准确地说，若 $a_1,a_2\in\ZZ$ 且 $b_1,b_2\in\ZZ$ 满足 $\overline{a_1}=\overline{b_1}$ 与 $\overline{a_2}=\overline{b_2}$，则
\[
  \overline{a_1+a_2}=\overline{b_1+b_2}
  \quad\text{且}\quad
  \overline{a_1a_2}=\overline{b_1b_2};
\]
也就是说，若
\[
  a_1\equiv b_1\pmod n\quad\text{且}\quad a_2\equiv b_2\pmod n,
\]
则
\[
  a_1+a_2\equiv b_1+b_2\pmod n
  \quad\text{且}\quad
  a_1a_2\equiv b_1b_2\pmod n.
\]
\noindent\textit{证明：} 假设 $a_1\equiv b_1\pmod n$，即 $a_1-b_1$ 可被 $n$ 整除。于是 $a_1=b_1+sn$，其中 $s$ 是某个整数。同样，$a_2\equiv b_2\pmod n$ 意味着 $a_2=b_2+tn$，其中 $t$ 是某个整数。因此
\[
  a_1+a_2=(b_1+b_2)+(s+t)n,
\]
所以 $a_1+a_2\equiv b_1+b_2\pmod n$，这说明剩余类的和与所选代表元无关。同样，
\[
  a_1a_2=(b_1+sn)(b_2+tn)=b_1b_2+(b_1t+b_2s+stn)n,
\]
这说明 $a_1a_2\equiv b_1b_2\pmod n$，因此剩余类的积也与所选代表元无关，证明完成。

我们稍后将看到，通过代表元相加来相加等价类的过程，是一个更一般构造（商构造）的特殊情形。等价类相加这个概念，在有理数加法的语境中已经很熟悉：每个有理数 $a/b$ 实际上都是一类表达式：$a/b=2a/2b=(-3a)/(-3b)$ 等等；我们经常改变代表元（例如通分）以便把两个分数相加（例如，计算 $1/2+1/3$ 时，改用与 $1/2$ 等价的代表元 $3/6$ 以及与 $1/3$ 等价的代表元 $2/6$，从而得到 $1/2+1/3=3/6+2/6=5/6$）。模算术的概念也很熟悉：在增加或减少若干小时后求一天中的时刻时，我们模 $12$ 化简并求最小剩余。

重要的是，要能够把某个等价关系的等价类看成可以操作的\emph{元素}（例如我们对分数所做的那样），而不是仅仅看成集合。与这种观点一致，我们常常把 $\ZZ/n\ZZ$ 的元素简单地记为 $\{0,1,\ldots,n-1\}$，其中加法和乘法都模 $n$ 化简。不过必须记住，$\ZZ/n\ZZ$ 的元素不是整数，而是通常整数的集合，并且其中的算术很不相同。例如，在整数 $\ZZ$ 中，$5+8$ 并不是 $1$，而在上面 $\ZZ/12\ZZ$ 的例子中却是如此。

能够在 $\ZZ/n\ZZ$ 中定义算术，在初等数论中有许多重要应用。作为一个简单例子，我们计算 $2^{1000}$ 的最后两位数字。首先注意，最后两位数字给出 $2^{1000}$ 除以 $100$ 的余数，因此我们关心的是包含 $2^{1000}$ 的模 $100$ 剩余类。计算得
\[
  2^{10}=1024\equiv24\pmod{100},
\]
于是
\[
  2^{20}=(2^{10})^2\equiv24^2=576\equiv76\pmod{100}.
\]
然后
\[
  2^{40}=(2^{20})^2\equiv76^2=5776\equiv76\pmod{100}.
\]
同样，$2^{80}\equiv2^{160}\equiv2^{320}\equiv2^{640}\equiv76\pmod{100}$。最后，
\[
  2^{1000}=2^{640}2^{320}2^{40}\equiv76\cdot76\cdot76\equiv76\pmod{100},
\]
所以最后两位数字是 $76$。

$\ZZ/n\ZZ$ 的一个重要子集由在 $\ZZ/n\ZZ$ 中有乘法逆元的剩余类构成：
\[
  (\ZZ/n\ZZ)^\times
  =\{\overline{a}\in\ZZ/n\ZZ\mid \text{存在 }\overline{c}\in\ZZ/n\ZZ\text{ 使得 }\overline{a}\cdot\overline{c}=\overline{1}\}.
\]
后面一些练习给出了一个证明大纲，说明 $(\ZZ/n\ZZ)^\times$ 也正是那些代表元与 $n$ 互素的剩余类全体，这证明了下面的命题。

\medskip
\noindent\textbf{命题 4.}
\[
  (\ZZ/n\ZZ)^\times=\{\overline{a}\in\ZZ/n\ZZ\mid (a,n)=1\}.
\]
容易看出，若 $\overline{a}$ 的任意一个代表元与 $n$ 互素，则所有代表元都与 $n$ 互素，因此命题右端的集合是良定的。

\paragraph{例}
当 $n=9$ 时，由命题可得
\[
  (\ZZ/9\ZZ)^\times=\{\overline{1},\overline{2},\overline{4},\overline{5},\overline{7},\overline{8}\}.
\]
这些元素的乘法逆元分别是
\[
  \{\overline{1},\overline{5},\overline{7},\overline{2},\overline{4},\overline{8}\}.
\]

若整数 $a$ 与 $n$ 互素，则欧几里得算法产生整数 $x$ 与 $y$，使得 $ax+ny=1$；于是 $ax\equiv1\pmod n$，所以 $\overline{x}$ 是 $\ZZ/n\ZZ$ 中 $\overline{a}$ 的乘法逆元。这给出了一个计算 $\ZZ/n\ZZ$ 中乘法逆元的有效方法。

\paragraph{例}
设 $n=60$ 且 $a=17$。应用欧几里得算法得到
\begin{align*}
  60&=(3)17+9,\\
  17&=(1)9+8,\\
  9&=(1)8+1.
\end{align*}
因此 $a$ 与 $n$ 互素，并且 $(-7)17+(2)60=1$。于是 $\overline{-7}=\overline{53}$ 是 $\ZZ/60\ZZ$ 中 $\overline{17}$ 的乘法逆元。

\subsection*{练习}

\begin{enumerate}[label=\arabic*.]
  \item 明确写出 $\ZZ/18\ZZ$ 的各个剩余类中的所有元素。
  \item 证明：$\ZZ/n\ZZ$ 中不同的等价类恰好是 $\overline{0},\overline{1},\overline{2},\ldots,\overline{n-1}$（使用带余除法）。
  \item 证明：若 $a=a_n10^n+a_{n-1}10^{n-1}+\cdots+a_1 10+a_0$ 是任意正整数，则
  \[
    a\equiv a_n+a_{n-1}+\cdots+a_1+a_0\pmod 9
  \]
  （注意这就是通常的算术规则：除以 $9$ 的余数与十进制各位数字之和模 $9$ 的余数相同；特别地，一个整数能被 $9$ 整除，当且仅当它的各位数字之和能被 $9$ 整除）[注意 $10\equiv1\pmod9$]。
  \item 计算 $37^{100}$ 除以 $29$ 的余数。
  \item 计算 $9^{1500}$ 的最后两位数字。
  \item 证明：$\ZZ/4\ZZ$ 中各元素的平方只有 $\overline{0}$ 和 $\overline{1}$。
  \item 证明：对任意整数 $a$ 与 $b$，$a^2+b^2$ 除以 $4$ 的余数永远不是 $3$（使用上一练习）。
  \item 证明：方程 $a^2+b^2=3c^2$ 在非零整数 $a,b,c$ 中没有解。[像前两个练习那样模 $4$ 考虑该方程，并证明 $a,b,c$ 都必须能被 $2$ 整除。于是 $a^2,b^2,c^2$ 各有因子 $4$，除以 $4$ 后说明原方程会有一组更小的解。不断迭代，得到矛盾。]
  \item 证明：任意奇整数的平方除以 $8$ 的余数总是 $1$。
  \item 证明：$(\ZZ/n\ZZ)^\times$ 的元素个数是 $\varphi(n)$，其中 $\varphi$ 表示欧拉 $\varphi$ 函数。
  \item 证明：若 $\overline{a},\overline{b}\in(\ZZ/n\ZZ)^\times$，则 $\overline{a}\cdot\overline{b}\in(\ZZ/n\ZZ)^\times$。
  \item 设 $n\in\ZZ$，$n>1$，且 $a\in\ZZ$ 满足 $1\leq a\leq n$。证明：若 $a$ 与 $n$ 不互素，则存在整数 $b$，满足 $1\leq b<n$ 且 $ab\equiv0\pmod n$；并推出不存在整数 $c$ 使得 $ac\equiv1\pmod n$。
  \item 设 $n\in\ZZ$，$n>1$，且 $a\in\ZZ$ 满足 $1\leq a\leq n$。证明：若 $a$ 与 $n$ 互素，则存在整数 $c$ 使得 $ac\equiv1\pmod n$。[使用两个整数的最大公因子是这两个整数的一个 $\ZZ$-线性组合这一事实。]
  \item 从前两个练习推出：$(\ZZ/n\ZZ)^\times$ 是 $\ZZ/n\ZZ$ 中满足 $(a,n)=1$ 的元素 $\overline{a}$ 的集合，并由此证明命题 4。在 $n=12$ 的情形中直接验证这一点。
  \item 对下列每一对整数 $a$ 与 $n$，证明 $a$ 与 $n$ 互素，并求出 $\ZZ/n\ZZ$ 中 $\overline{a}$ 的乘法逆元。
  \begin{enumerate}[label=(\alph*)]
    \item $a=13,\ n=20$。
    \item $a=69,\ n=89$。
    \item $a=1891,\ n=3797$。
    \item $a=6003722857,\ n=77695236973$。[对这两个整数，欧几里得算法只需要 $3$ 步。]
  \end{enumerate}
  \item 编写一个计算机程序，对任意给定的 $n$ 作模 $n$ 的加法和乘法。这些运算的输出应为两个整数的和与积的最小剩余。同时加入如下功能：若 $(a,n)=1$，则在需要时输出一个介于 $1$ 与 $n-1$ 之间的整数 $c$，使得 $\overline{a}\cdot\overline{c}=\overline{1}$。（当然，你的程序不应直接调用许多系统中已经内置的 ``mod'' 函数。）
\end{enumerate}
